%% =====================================================================
%%  B1-T-01-01 -- what B1 contributes to "The Five-Percent Premise"
%%  -------------------------------------------------------------------
%%  LAYER A: APPEND-ONLY. Written once, then left alone. Material found
%%  only here sits in the `unique` slot and is protected by rule R10.
%% =====================================================================
%% @kind: source
%% @id: B1-T-01-01
%% @book: B1
%% @topic: T-01-01
%% @locator: ch. 1, pp. 1--8
%% @status: distilled
%% @unique: yes
%% @integrated: yes
%% @updated: 2026-09-19

\dossier{B1}{T-01-01}{\bookshort{B1}}{ch. 1, pp. 1--8}

\begin{distilled}
  \item Framing claim of the book: about five percent of people manipulate the other ninety-five percent, and the difference is knowledge of human nature rather than intelligence.
  \item The author's stated method is to learn from practitioners --- hustlers and con artists --- rather than from academic psychology, on the ground that a tactic is only real if it works in the field.
  \item The chapter closes by turning the premise defensive: the point of learning the tactics is to stop being on the receiving end of them.
\end{distilled}

\begin{unique}
  \item The five-percent figure itself, and the framing of manipulation as a learnable trade acquired from practitioners rather than from theory.
\end{unique}

\begin{terms}
  \item \emph{street-educated mind} --- Sparkman's term for judgment acquired from repeated losses rather than from study.
\end{terms}
